\section{Related Work}
\label{sec:related-work}

\subsection{Production L2 State Commitment Patterns}

The three major production ZK rollups---zkSync Era, Polygon zkEVM, and Scroll---each
implement distinct L1 state commitment patterns that inform our design space.

\textbf{zkSync Era} uses a three-phase commit-prove-execute lifecycle~\cite{zksync2024docs,
quarkslab2023zksync}. The \texttt{commitBatches} transaction posts batch metadata and
pubdata hashes; \texttt{proveBatches} submits the ZK-SNARK proof to a specialized
\texttt{Verifier} contract; and \texttt{executeBatches} finalizes state by processing
L2$\to$L1 messages and updating the global state root. This separation enables batch
aggregation---multiple committed batches can be proved together---at the cost of 2--3
storage slots (64--96 bytes) per batch and three separate L1 transactions per batch
lifecycle. The state root is a state-diff commitment: zkSync publishes only changed
storage slots, not full transaction data.

\textbf{Polygon zkEVM} uses a two-phase sequence-verify model~\cite{polygon2024zkevm,
polygonzkevm2024contract}. The trusted sequencer calls \texttt{sequenceBatches} to post
transaction data and an accumulated input hash (\texttt{accInputHash}) that chains batches
cryptographically: $h_i = H(h_{i-1}, \text{txHash}, \text{exitRoot}, \text{timestamp},
\text{seqAddr})$. The aggregator then calls \texttt{verifyBatches} with a SNARK proof
covering one or more sequenced batches. Per-batch storage requires approximately 2 slots
(64 bytes): the \texttt{SequencedBatchData} struct and batch counter.

\textbf{Scroll} uses a two-phase commit-finalize pattern~\cite{scroll2024docs,
scrollchain2024contract}. The \texttt{commitBatchWithBlobProof} transaction posts a batch
header hash plus EIP-4844 blob data~\cite{eip4844}. The \texttt{finalizeBatchWithProof}
transaction submits a PLONK proof along with pre- and post-state roots. Per-batch storage
requires 3 slots (96 bytes): batch hash, finalized state root, and withdraw root. Scroll's
V7 batch format achieves approximately 90\% reduction in commitment costs by moving most
metadata to blobs.

\subsection{Comparison of Production Storage Patterns}

Table~\ref{tab:production-comparison} summarizes the storage patterns across production
systems and positions our target design.

\begin{table}[htbp]
\centering
\caption{Production L2 State Commitment Comparison}
\label{tab:production-comparison}
\small
\begin{tabular}{@{}lcccc@{}}
\toprule
\textbf{System} & \textbf{Phases} & \textbf{Slots} & \textbf{Bytes} & \textbf{Chaining} \\
\midrule
zkSync Era       & 3 & 2--3 & 64--96 & Batch hash \\
Polygon zkEVM    & 2 & 2    & 64     & accInputHash \\
Scroll           & 2 & 3    & 96     & parentBatchHash \\
\midrule
\textbf{Basis (ours)} & \textbf{1} & \textbf{1} & \textbf{32} & \textbf{prevRoot chain} \\
\bottomrule
\end{tabular}
\end{table}

A key distinction is that all production systems separate commitment and verification into
multiple L1 transactions, while our design executes both atomically in a single transaction.
This is feasible in the enterprise validium model because: (1) the operator is a known,
authorized enterprise rather than an untrusted sequencer; (2) there is no need for
challenge periods or fraud proofs; and (3) proof generation and submission are tightly
coupled within the enterprise node.

\subsection{EVM Gas Mechanics}

The EVM gas schedule, as modified by EIP-2929~\cite{eip2929} and EIP-2200~\cite{eip2200},
defines the cost structure for on-chain state management. Cold \texttt{SSTORE} operations
(first access to a slot within a transaction) cost 22,100 gas for zero-to-nonzero writes
and 5,000 gas for nonzero-to-nonzero updates. Warm accesses reduce these to 20,000 and
2,900 gas respectively. These costs are identical on Avalanche Subnet-EVM, which inherits
the full go-ethereum gas schedule~\cite{avalanche2024subnetdocs}.

For Groth16 verification, the BN256 precompile gas costs were reduced by EIP-1108~\cite{eip1108}
to 150 gas for point addition, 6,000 gas for scalar multiplication, and $45{,}000 + 34{,}000
\cdot k$ gas for pairing with $k$ pairs. A Groth16 proof with $N$ public inputs requires
$N$ scalar multiplications, $N$ additions, and one pairing check with 4 pairs, yielding a
verification cost of $N \cdot 6{,}150 + 181{,}000$ gas~\cite{groth2016size}.

\subsection{Enterprise ZK Systems}

ZoKrates~\cite{eberhardt2018private} pioneered practical ZK proof generation for Ethereum
smart contracts but targets single-proof use cases rather than continuous batch submission
pipelines. Hyperledger Fabric~\cite{hyperledger2024fabric} provides enterprise privacy
through channel isolation but without validity proofs---participants must trust the ordering
service. StarkWare's validium concept~\cite{starkware2024validium} introduced the pattern
of on-chain validity proofs with off-chain data availability, which we adapt for
per-enterprise state isolation on an Avalanche L1.

No prior work has analyzed the gas cost decomposition of combined ZK verification and state
commitment in a single atomic transaction for enterprise validium, nor has the specific
constraint that pairing verification imposes on storage layout budget been quantified.
